\section{Capítulo~\ref{sec:sensors}. As entradas da máquina}

A construção dos sensores foi medida diretamente, com registros da corrente de células isoladas, das vibrações da cóclea e com contagens de células da retina; o limite de sensibilidade e a fórmula do ruído balístico decorrem da física, e que os códigos dos sensores estejam ajustados para a máxima informação é uma hipótese com apoio forte.

{\footnotesize
\setlength{\tabcolsep}{3pt}\renewcommand{\arraystretch}{1.15}
\begin{longtable}{>{\raggedright}p{0.5cm}>{\raggedright}p{4.0cm}>{\raggedright}p{5.6cm}>{\raggedright}p{2.3cm}>{\raggedright\arraybackslash}p{1.7cm}}
\toprule
N.º & Proposição & Experimento e o que foi medido & Fonte & Status \\
\midrule\endfirsthead
\toprule
N.º & Proposição & Experimento e o que foi medido & Fonte & Status \\
\midrule\endhead
1 & O sensor é um transdutor: o receptor gera corrente, e a saída das células sensoriais do olho e do ouvido é glutamato para o barramento \nt{GLUT}; as primeiras células não emitem disparos (fig.~\ref{fig:transducer}) & Bastonetes e cones de tartaruga liberam um aminoácido excitatório: ele é captado por canais de glutamato num fragmento de membrana excisado. Células ciliadas internas da cóclea de rato: as correntes nas terminações das fibras nervosas passam por receptores AMPA de glutamato, e a taxa de liberação cresce com a despolarização gradual da célula até $150$\,Hz & \cite{CopenhagenJahr1989,GlowatzkiFuchs2002} & medido \\
2 & O sensor de luz no escuro responde a um único fóton com uma corrente de cerca de um picoampere & Eletrodo de sucção no segmento externo de bastonete de sapo: as respostas a flashes fracos são quantizadas, evento de $\approx1$\,pA ($5\%$ da corrente de escuro) com dispersão de $0.2$\,pA, eficiência $0.5$. Humanos relatam a chegada de um único fóton acima do acaso & \cite{Baylor1979,Tinsley2016} & medido \\
3 & O sensor de som percebe um deslocamento menor que um nanômetro & Método Mössbauer, membrana basilar da cóclea de porquinho-da-índia no ponto de $16$--$19$\,kHz: no limiar de resposta do nervo auditivo, a velocidade da membrana é de cerca de $0.04$\,mm/s, isto é, um deslocamento de $\approx0.35$\,nm (cálculo nosso). Mediu-se a membrana, e não o feixe do sensor & \cite{Sellick1982,Hudspeth1989} & medido \\
4 & No escuro, a precisão é limitada pelo ruído balístico dos fótons~\eqref{eq:photonnoise}; com luz fraca, a acumulação é mais longa & A fórmula decorre da estatística de Poisson. No bastonete, o ruído com luz fraca coincide com a soma de eventos quânticos independentes; sobre um fundo claro, a resposta ao flash é menor e mais curta. O valor ``décimos de segundo'' não foi verificado separadamente aqui & \cite{Baylor1979} & teoria \\
5 & A faixa de entrada é de dez ordens de grandeza para a luz e doze para o som, e a taxa de disparos, menos de três ordens & Para o som: a resposta da membrana basilar na frequência característica cresce com inclinação de até $0.2$\,dB/dB na faixa de $40$--$80$\,dB, e a compressão aparece também acima de $100$\,dB. Os próprios números de ordens são estimativas padrão, sem fonte no capítulo & \cite{Ruggero1997} & medido \\
6 & A resposta do sensor segue~\eqref{eq:nakarushton}, e $\sigma$ acompanha o brilho médio, deslocando a curva ao longo do eixo logarítmico (fig.~\ref{fig:adaptation}) & A fórmula foi ajustada primeiro aos potenciais S da retina de peixe. Cones de tartaruga: as respostas seguem $V=I V_m/(I+\sigma)$ e cobrem $\sim3.5$ ordens de brilho; após $2$--$3$ minutos de fundo, a curva se desloca na horizontal, mantendo a largura. A ligação com a divisão pela atividade total é de uma revisão & \cite{NakaRushton1966,NormannPerlman1979,Carandini2012} & medido \\
7 & Os sensores transmitem mudanças, e os seus códigos estão ajustados para a máxima informação por disparo (proposição~\ref{prop:sensors}) & O princípio foi proposto teoricamente. A evidência mais forte: na mosca, a característica ``contraste $\to$ resposta'' dos neurônios da primeira camada coincide com a distribuição acumulada de contrastes das cenas naturais, e assim se obtém a máxima informação & \cite{Barlow1961,Laughlin1981} & hipótese \\
8 & Os sensores de ligar e de desligar formam um código de dois fios; a câmera de eventos o repete em silício & Revisão de experimentos: os canais de ligar e de desligar da retina se separam já na primeira sinapse e podem ser silenciados separadamente. Câmera: $128\times128$ pixels sem quadros, faixa de $120$\,dB, latência de $15$\,\textmu s & \cite{Schiller1992,Lichtsteiner2008,Gallego2022} & medido \\
9 & O olho tem $\sim10^8$ sensores de luz e $\sim10^6$ neurônios de saída: compressão de $\sim100$ vezes & Contagem em preparações inteiras de retina humana: em média $92$ milhões de bastonetes e $4.6$ milhões de cones; células de saída (ganglionares), de $0.7$ a $1.5$ milhão conforme o olho & \cite{Curcio1990,CurcioAllen1990} & medido \\
10 & O camundongo tem mais de trinta canais de saída do olho & Camundongo: imageamento de cálcio de dois fótons de mais de $11\,000$ células de saída e agrupamento das respostas: bem mais de $30$ tipos funcionais. Para humanos, o número não foi medido & \cite{Baden2016} & medido \\
11 & O olho do porquinho-da-índia transmite $\sim875\,000$ bit/s; o recálculo para o olho humano dá $\sim10^7$ bit/s & Porquinho-da-índia, matriz de múltiplos eletrodos, movimento em cenas naturais: $6$--$13$ bit/s por célula, $\sim875\,000$ bit/s para $\sim10^5$ células de saída. Para humanos ($\sim10^6$ células), $\sim10^7$ bit/s é um recálculo dos autores & \cite{Koch2006} & medido \\
12 & O ouvido é um banco de filtros passa-faixa com um mapa de frequências quase logarítmico (fig.~\ref{fig:filterbank}) & O ponto de máxima vibração da membrana basilar depende da frequência; a função ``frequência--lugar'' é quase exponencial e descreve com uma só forma as cócleas humanas e de animais vivos & \cite{Greenwood1990,Robles2001} & medido \\
13 & Os ressonadores são ativos: parte dos sensores amplifica as vibrações fracas milhares de vezes em amplitude ($66$--$76$\,dB), e o ganho cai com o volume & Vibrometria a laser da base da cóclea de chinchila: com tons fracos na frequência característica, ganho de $66$--$76$\,dB em relação ao estribo; a morte reduz a sensibilidade em $60$--$81$\,dB ($10^3$--$10^4$ vezes em amplitude) e elimina a compressão. A fonte da força são as células ciliadas externas & \cite{Ruggero1997,Robles2001} & medido \\
14 & O amplificador do ouvido está no limite da autoexcitação & Cálculo: perto de um ponto de bifurcação de Hopf são inevitáveis a compressão da faixa, a sintonia aguda e os tons de combinação, isto é, todas as não linearidades conhecidas da audição. Que a cóclea esteja ajustada assim não foi provado diretamente & \cite{Eguiluz2000} & hipótese \\
15 & Nas frequências baixas, os disparos vêm em fase com o som (no porquinho-da-índia, o travamento enfraquece a partir de $\sim600$\,Hz e é perceptível até $\sim3.5$\,kHz), e por eles se acha a direção; nas altas, resta o código de lugar & Nervo auditivo de porquinho-da-índia: o travamento de fase começa a cair em torno de $600$\,Hz e desaparece acima de $3.5$\,kHz; no gato e no macaco o limite fica cerca de uma oitava acima. Diferença de tempo de chegada aos dois ouvidos: revisão & \cite{PalmerRussell1986,Grothe2010} & medido \\
16 & Os demais sensores (tab.~\ref{tab:sensors}): o olfato tem $\sim400$ tipos de receptores, um por neurônio, e código combinatório; o paladar funciona em parte por ATP e serotonina; o tato tem sensores de adaptação rápida e lenta; calor e frio são separados & Camundongo: um receptor reconhece muitos cheiros, um cheiro, muitos receptores, e cheiros diferentes, conjuntos diferentes. Sem receptores de ATP, o nervo do paladar não responde; os botões gustativos liberam serotonina. Registro de fibras isoladas da pele da mão humana: quatro tipos pela velocidade de adaptação. O canal do frio é distinto dos canais do calor & \cite{Mombaerts2004,Malnic1999,Finger2005,Huang2005,JohanssonVallbo1979,McKemy2002} & medido \\
\bottomrule
\end{longtable}}
